\begin{definition}
Let \(D\) be a \noderef{HierarchySituationTwoTailData} for the fixed
minimum counterexample, and put
\[
 B=\noderef{HierarchyFullTailPrefix}(C,\ell-2),\qquad
 B^+=\noderef{HierarchyFullTailPrefix}(C,\ell-1),\qquad
 \ell=\noderef{HierarchyFullTailLength}(C).
\]
A pair of \emph{Situation 7.2 reordered trees} consists of an exact
index \(r\in\operatorname{Fin}|T|\) with \(r+2=|T|\), the individual
edge \(e_{p-1}=T.\mathit{edge}(r)\), and two ordered tree-sequences
\(T^-\) and \(T^*\) satisfying all of the following.

The original display is obtained by the two exact one-edge extensions
\[
 B\mathbin{+}(e_{p-1},y_{p-1})=B^+,\qquad
 B^+\mathbin{+}(e_p,y_p)=T.
\]
The reordered displays are obtained by
\[
 B\mathbin{+}(e_p,y_p)=T^-,\qquad
 T^-\mathbin{+}(e_{p-1},y_{p-1})=T^*.
\]
Each plus sign here is the exact ordered
\noderef{TreeSequenceAdjoinEdge} relation, and the displayed last
vertices are required literally.  Moreover
\[
 \begin{aligned}
 V(T^-)&=V(T)\setminus\{y_{p-1}\},&
 E(T^-)&=E(T)\setminus\{e_{p-1}\},\\
 V(T^*)&=V(T),& E(T^*)&=E(T).
 \end{aligned}
\]
Thus \(T^-\) is the manuscript's tree obtained by deleting
\(y_{p-1}\) together with its displayed attachment edge, while \(T^*\)
interchanges the final two attachment pairs.  The equalities concern
sets of individual edge identities, not endpoint pairs.
\end{definition}
